\mysection{OpenMP}
\label{openmp}
\myindex{OpenMP}

OpenMPは、単純なアルゴリズムを並列化する最も簡単な方法の一つです。

\myindex{Bitcoin}

例として、暗号学的\IT{nonce}を計算するプログラムを構築してみましょう。

私の簡単な例では、
\IT{nonce}は特定の特徴を持つハッシュを生成するために、平文の暗号化されていないテキストに追加される数値です。

例えば、Bitcoinプロトコルのある段階では、結果のハッシュが特定の数の連続するゼロを含むような\IT{nonce}を見つけることが要求されます。
これは\q{proof of work}とも呼ばれます
\footnote{\href{http://go.yurichev.com/17101}{wikipedia}}
（つまり、システムが集約的な計算を行い、そのために時間を費やしたことを証明します）。

\myindex{SHA512}
私の例はBitcoinとは何の関係もありません。
\q{hello, world!\_}文字列に数値を追加して、
\q{hello, world!\_<number>}がSHA512アルゴリズムでハッシュ化されたときに、少なくとも3つのゼロバイトを含むような数値を見つけようとします。

ブルートフォースを0..INT32\_MAX-1の間隔に制限します（つまり、\TT{0x7FFFFFFE}または2147483646）。

アルゴリズムは非常に簡単です：

\lstinputlisting[style=customc]{advanced/700_openmp/openmp_example.c}

\TT{check\_nonce()}関数は、文字列に数値を追加し、
SHA512アルゴリズムでハッシュ化し、結果内の3つのゼロバイトをチェックします。

コードの非常に重要な部分は次のとおりです：

\begin{lstlisting}[style=customc]
	#pragma omp parallel for
	for (i=0; i<INT32_MAX; i++)
		check_nonce (i);
\end{lstlisting}

はい、それだけ簡単です。\TT{\#pragma}なしでは、
0から\TT{INT32\_MAX}（\TT{0x7fffffff}または2147483647）までの各数値に対して\TT{check\_nonce()}を呼び出すだけです。
\TT{\#pragma}があると、コンパイラはループ間隔をより小さなものに分割し、
利用可能なすべての\ac{CPU}コア上で実行するための特別なコードを追加します
\footnote{注意：これは意図的に最も簡単な例ですが、実際には、OpenMPの使用はより困難で複雑になる可能性があります}。

例をコンパイルできます
\footnote{sha512.(c|h)とu64.hファイルはOpenSSLライブラリから取得できます：
\url{http://go.yurichev.com/17324}}
MSVC 2012で：
% FIXME1: \footnote{other source code files can be downloaded here: ...} 

\begin{lstlisting}
cl openmp_example.c sha512.obj /openmp /O1 /Zi /Faopenmp_example.asm
\end{lstlisting}

またはGCCで：

\begin{lstlisting}
gcc -fopenmp 2.c sha512.c -S -masm=intel
\end{lstlisting}

\subsection{MSVC}

MSVC 2012がメインループを生成する方法は次のとおりです：

\lstinputlisting[caption=MSVC 2012,style=customasmx86]{advanced/700_openmp/MSVC_loop.asm}

\TT{vcomp}で始まるすべての関数は
OpenMP関連で、vcomp*.dllファイルに格納されています。
ここでスレッドのグループが開始されます。

\TT{\_main\$omp\$1}を見てみましょう：

\lstinputlisting[caption=MSVC 2012,style=customasmx86]{advanced/700_openmp/MSVC_1.asm}

この関数は$n$回並列で開始されます。ここで$n$は\ac{CPU}コア数です。\\
\TT{vcomp\_for\_static\_simple\_init()}は
現在のスレッド番号に応じて、現在のスレッドのfor()構造の間隔を計算します。

ループの開始値と終了値は\TT{\$T1}と\TT{\$T2}ローカル変数に格納されます。
\TT{vcomp\_for\_static\_simple\_init()}関数への引数として\TT{7ffffffeh}（または2147483646）も見ることができます
---これは、均等に分割される全体のループの繰り返し数です。

その後、\TT{check\_nonce()}関数への呼び出しがある新しいループが見えます。これがすべての作業を行います。

統計を収集するために、\TT{check\_nonce()}関数の呼び出しについて、関数が呼び出された引数について
\TT{check\_nonce()}関数の最初にコードを追加してみましょう。

実行すると次のようになります：

\begin{lstlisting}
threads=4
...
checked=2800000
checked=3000000
checked=3200000
checked=3300000
found (thread 3): [hello, world!_1611446522]. seconds spent=3
__min[0]=0x00000000 __max[0]=0x1fffffff
__min[1]=0x20000000 __max[1]=0x3fffffff
__min[2]=0x40000000 __max[2]=0x5fffffff
__min[3]=0x60000000 __max[3]=0x7ffffffe
\end{lstlisting}

はい、結果は正しく、最初の3バイトはゼロです：

\begin{lstlisting}
C:\...\sha512sum test
000000f4a8fac5a4ed38794da4c1e39f54279ad5d9bb3c5465cdf57adaf60403
df6e3fe6019f5764fc9975e505a7395fed780fee50eb38dd4c0279cb114672e2 *test
\end{lstlisting}

実行時間は4コアIntel Xeon E3-1220 3.10 GHzで約2..3秒です。
タスクマネージャでは5つのスレッドが見えます：
1つのメインスレッド + 4つの追加。
この例を小さく明確に保つために、さらなる最適化は行われていません。
しかし、おそらくもっと速くできるでしょう。
私の\ac{CPU}は4コアを持っているため、OpenMPは正確に4つのスレッドを開始しました。

統計表を見ることで、ループが4つの均等な部分にどのように分割されたかが明確にわかります。
ああ、まあ、最後のビットを考慮しなければほぼ均等です。

\glslink{atomic operation}{atomic operations}のためのプラグマもあります。

このコードがどのようにコンパイルされるか見てみましょう：

\begin{lstlisting}[style=customc]
	#pragma omp atomic
	checked++;

	#pragma omp critical
	if ((checked % 100000)==0)
		printf ("checked=%d\n", checked);
\end{lstlisting}

\lstinputlisting[caption=MSVC 2012,style=customasmx86]{advanced/700_openmp/MSVC_2.asm}

結局のところ、vcomp*.dllの\TT{vcomp\_atomic\_add\_i4()}関数は、
\TT{LOCK XADD}命令\footnote{
LOCKプレフィックスについて詳しく読む：\myref{x86_lock}}を持つ小さな関数です。

\TT{vcomp\_enter\_critsect()}は
最終的にwin32 \ac{API}関数\\
\TT{EnterCriticalSection()}を呼び出します
\footnote{クリティカルセクションについて詳しく読める場所：\myref{critical_sections}}。

\subsection{GCC}

GCC 4.8.1は
まったく同じ統計表を表示するプログラムを生成します。

そのため、GCCの実装は同じ方法でループを部分に分割します。

\begin{lstlisting}[caption=GCC 4.8.1,style=customasmx86]
	mov	edi, OFFSET FLAT:main._omp_fn.0
	call	GOMP_parallel_start
	mov	edi, 0
	call	main._omp_fn.0
	call	GOMP_parallel_end
\end{lstlisting}

MSVCの実装とは異なり、GCCのコードは3つのスレッドを開始し、
4番目を現在のスレッドで実行します。そのため、MSVCの5つではなく4つのスレッドです。

\TT{main.\_omp\_fn.0}関数は次のとおりです：
 
\lstinputlisting[caption=GCC 4.8.1,style=customasmx86]{advanced/700_openmp/GCC_1.asm}

ここで分割が明確に見えます：\TT{omp\_get\_num\_threads()}と\TT{omp\_get\_thread\_num()}を呼び出すことで

実行中のスレッド数と現在のスレッド番号を取得し、その後ループの間隔を決定します。
その後、\TT{check\_nonce()}を実行します。

GCCは、別のDLL関数への呼び出しを生成するMSVCとは異なり、
\TT{LOCK ADD}命令をコードに直接挿入しました：

\lstinputlisting[caption=GCC 4.8.1,style=customasmx86]{advanced/700_openmp/GCC_2.asm}

GOMPプレフィックスを持つ関数はGNU OpenMPライブラリからのものです。
vcomp*.dllとは異なり、そのソースコードは自由に利用できます：
\href{http://go.yurichev.com/17102}{GitHub}。